\section{Discussion}
\label{sec:discussion}

\subsection{Attention is common and conversion is not}

The evidence sits at three levels. On attention the result is strong and shared by all four systems. They refuse, reprice or publicly contest trades by reference to what the trade gives a rival, and their notes say so. On pricing the result is directional. Asks to completing buyers are higher in every coarse cut, but the matched cases are mixed, bids show no common direction, and the stated arithmetic is sometimes wrong. On outcomes, attention is not what separates the systems. Grok refuses \fact{d.holder.Grok.300plus.rejectpct}\% of completion requests while it is liquid and states a rival consequence in \fact{inh.reason.D.both.Grok} of its \fact{inh.reason.rows.Grok} refusals of its own completion, yet it built at once on \fact{inh.window.Grok.built} of \fact{inh.window.Grok.groups} groups and collected \$\fact{d.conv.rpl.Grok} per opponent landing. Sol counters where Grok refuses, sells a blocker when it has something to build, and built at once on \fact{inh.window.Sol.built} of \fact{inh.window.Sol.groups}. Gemini refuses as often as Grok and builds like Sol. On question A, then, this cohort points one way. Speaking and acting on a rival-aware view is common to these systems. Turning a position into income is rarer, and it is where their trajectories differ.

None of this shows that building caused winning. Opus's building gap is partly its response condition, and Grok's is partly its conflicting calls. Wins of \fact{d.wins.Sol}, \fact{d.wins.Gemini}, \fact{d.wins.Opus} and \fact{d.wins.Grok} by system sit beside wins of \fact{conv.cohort.wins-by-seat.1}, \fact{conv.cohort.wins-by-seat.2}, \fact{conv.cohort.wins-by-seat.3} and \fact{conv.cohort.wins-by-seat.4} by seat, and twenty games cannot tell skill from seat and dice.

\subsection{Two results from bargaining theory}

The observations fit two established results. Jehiel and Moldovanu show that negative externalities among potential buyers can delay a sale under complete information, because what a party will pay depends on who would otherwise obtain the object~\cite{jehiel1995delay}. Our holders are in a stronger version of that position, since the seller itself will pay rent to the buyer. They also tie their result to Rubinstein's point that the cost of rejecting an offer determines the outcome~\cite{rubinstein1982perfect}, and a holder short of cash has a high cost of waiting. Shleifer and Vishny show that a distressed seller receives less than an asset's value in best use when its natural buyers are constrained~\cite{shleifer1992liquidation,shleifer2011firesales}. Here the natural buyer of a completing deed is a single rival. Even so, three of the four systems mostly sold above list price when short of cash, and the low prices came from Opus. We test neither theory. They suggest that this behavior is what bargaining under externalities and cash constraints tends to produce, and that it is not peculiar to these systems.

\subsection{What would count against this account}

Each claim has an observation that would contradict it. The liquidity account fails if, in a new cohort, liquid holders release completing deeds as readily as cash-poor ones, or if holders who sell after refusing have not lost cash in between. The pricing claim reduces to composition if asks stop rising with the buyer's proximity once deed and state are matched in adequate numbers. With seven matched cases, that claim is the weakest here. The conversion contrast is a cohort accident if sellers who build after selling win no more often than those who do not. The reliability of forced finance fails if payable debts go unpaid at a material rate under the same menu.

\subsection{The motivating questions}

Table~\ref{tab:questions} states what this cohort shows about each of the eight questions and A to C, and what it does not measure. An unmeasured construct is not a negative finding about capability.

\begingroup\footnotesize
\renewcommand{\arraystretch}{1.12}
\begin{longtable}{@{}L{0.20\linewidth}L{0.45\linewidth}L{0.29\linewidth}@{}}
\caption{Coverage of the motivating questions. ``Inherited'' marks results from earlier work that this paper did not rerun.}\label{tab:questions}\\
\toprule
Question & What this cohort shows & Not measured\\
\midrule\endfirsthead
\toprule
Question & What this cohort shows & Not measured\\
\midrule\endhead
\bottomrule\endfoot
1. Strategic reasoning & Agents routinely reject trades that would help them because of what a rival gains. They reject \fact{barg.census.cells.recipient-only.reject} of \fact{barg.census.cells.recipient-only.n} offers completing only their own group and state a rival consequence in \fact{inh.reason.D.both} of \fact{inh.reason.rows} such refusals (inherited readings). & Whether any refusal was correct. Regret and exploitability.\\
2. Opponent modeling & Agents track rivals' public cash and building capacity, send completion requests to holders short of cash (\fact{d.target.completion.pct}\% against \fact{d.target.other.pct}\% for other requests) and name who would inherit an estate. & Accuracy or calibration of beliefs about hidden intentions.\\
3. Negotiation, cooperation and deception & A full response census. Counters reprice the same bundle. Incidents show authorized bluffs, selective disclosure, a withdrawn assurance and two public appeals to deny a rival deeds or funds. Acceptance is \fact{barg.census.responses-all.accept} of \fact{barg.census.responses-all.n} responses, so no general bias toward agreement appears. & Persuasion effects, joint surplus, fairness and the prevalence of deception.\\
4. Planning and memory & Stated timing rules that match later population behavior (Sol waits for opponents to need cash). Systems differ in whether cash is held or built with. One multi-turn building target was reached in G9 (inherited). & Benefit of memory and plan-completion rates. Opus mostly lacked its own notes.\\
5. Adaptation and learning & Holders change terms as their cash changes. After a corrective message, 9 of 50 later build submissions repeated an uneven-build error, against 54 of 59 without one, all of the latter Opus (inherited). & Learning caused by feedback. Adaptation to unfamiliar opponents.\\
6. Economic decisions and risk & Prices relative to list by role and position. \fact{debt.paid} of \fact{debt.feasible-episodes} payable debts paid. Rent per landing and construction after sales by system. & Whether any price or reserve was right. Risk-adjusted returns.\\
7. Collective dynamics & Estates concentrate deeds with eventual winners (\fact{flows.estate-deeds-nonterminal-to-winner} of \fact{flows.estate-deeds-nonterminal} before the last elimination). Estates end \fact{barg.dyads.unsettled-resolved-by-estate} of \fact{barg.dyads.unsettled} unsettled stand-offs. Appeals to deny a rival drew partial follow-through. & Coalition stability and the effect of any appeal.\\
8. Validity and generalization & Event-derived state matches all \fact{track.prompts} prompts. Conditions are stated with the results they touch, and main contrasts are checked by seed block. & A ranking of models, play under a repaired interface, other rosters.\\
\midrule
A. Strategic or sounding strategic & Rival-aware talk and action are common. Conversion into income differs by system. One note finds a payment route that the applied action abandons (Incident~\ref{inc:g12}). & A general relation between the quality of explanations and the quality of play.\\
B. General and multi-agent capability & Nothing. The cohort holds no measure of general capability. & Any transfer relationship.\\
C. Skill, luck and weak opponents & Win counts by system and by seat are similarly uneven. The first creditor already led in \fact{estate.first-creditor-leads-assets} of 20 games, and one payable bankruptcy handed a last-place player a developed estate. & The shares of skill and chance, and any counterfactual winner.\\
\end{longtable}
\endgroup

\section{Limitations}
\label{sec:limits}

The design is small and partly confounded. It has five seeds, one roster and one cyclic seating order, and opponents see each other's model names. Opus's response condition coincides with the model and shaped the games its opponents played. Grok's applied calls sometimes contradict its text. The agents were not told the rules, and most prompts omit prices, so a wrong figure in a note can be a recall error as easily as a reasoning error.

The comparisons are descriptive. Offers, pairs and games are dependent, and we report block-level consistency in place of tests. The cash thresholds of \$100 and \$300, the eight-turn building window and the asset valuation in Figure~\ref{fig:estates} were chosen for this analysis after earlier work on the same records, and we give alternatives where they matter. No decision is scored against a value reference, so we cannot say that any refusal, price or sale was right. Seller and buyer outcomes are conditioned on positions that the comparisons do not hold fixed.

Several results are inherited and were not rerun here, namely the payment sequences, the first-management-turn classification, the full replay and the semantic readings. The readings are model-assisted, come from one model family and have no human check. A research assistant chose the eight incidents after reading the records, to show distinct mechanisms. Grok appears in seven of them. Its private notes are the longest, at a median of \fact{cond.channels.Grok.median-note-characters} characters against \fact{cond.channels.Sol.median-note-characters} for Sol and \fact{cond.channels.Gemini.median-note-characters} for Gemini, and only its responses regularly carry text beside the call. The incidents estimate no rate.

\section{Related work}

\paragraph{LLM negotiation.} NegotiationArena studies bilateral scenarios with defined payoffs and reports that an agent feigning desperation raises its payoff by 20\% against one opponent model~\cite{bianchi2024negotiationarena}. Incident~\ref{inc:g13} is a distress pitch of that kind in a four-player game with no payoff measure. Cosentino and colleagues give agents a ranked summary of the counterparty's preferences in a controlled multi-attribute bargain. Their agents model the counterparty accurately and do not turn that into better agreements, and final prices track opening anchors~\cite{cosentino2026counterparty}. Our contrast between attention and conversion is the same distinction observed in play, without their control or their utility measure. Earlier work built scorable multi-party negotiation games with greedy and adversarial players~\cite{abdelnabi2024negotiation} and trained end-to-end negotiators on a bargaining task with hidden rewards~\cite{lewis2017deal}.

\paragraph{Long-horizon economic agents.} Li and colleagues analyze 2,583 emails between agents in twenty year-long Vending-Bench Arena runs. They label misaligned communication with a frozen codebook, check it against a hand-labeled reference set and other judge models, and estimate uncertainty at the run level. They report that low inventory raises the odds of a misaligned email~\cite{li2026arena}. Their prompts carry no conduct instruction, while ours invites deception, so their misalignment framing does not transfer to our records. Their finding that economic pressure changes communication parallels ours that cash pressure changes bargaining, and their label validation is stronger than ours.

\paragraph{Monopoly and board-game benchmarks.} AgentMonopoly compares model families, sizes and personas through trade proposals, acceptance rates, rent waivers and cash management, in an environment that eliminates a player when its cash goes negative~\cite{wu2026agentmonopoly}. FinBoardBench evaluates nine models on three financial board games, including a Monopoly with accept, reject and counter responses, and reports that models buy assets when short of cash and undervalue liquidity~\cite{hu2026finboardbench}. Our records add a liquidation menu, intact transfer of developed estates and a complete bargaining history. The liquidity failure they show comes after acquisition, in what holders do with a completed group or a sale's proceeds. Cattle Trade finds that two rule-based agents beat most of the LLMs it tests and that spending discipline tracks rank more closely than any single skill~\cite{muller2026cattletrade}. That agrees with our conversion finding and is a reminder that our cohort has no scripted baseline.

\section{Conclusion}

In twenty four-player Monopoly games, four LLM systems all bargained as if a deed were a rival's claim on a monopoly. They refused to sell such deeds, asked more for them and argued in public about who should fund or inherit them. The deeds moved when their holders ran short of cash. What then separated the systems was whether cash and completed groups became houses. Paying a payable debt, a task with an explicit menu and an immediate check, was done reliably by all four. The cohort is small, observational and uneven in its response conditions, and it ranks no model. It identifies liquidity management after a strategic position is won as the place where these agents' play diverges, and it gives four concrete observations that a new cohort could contradict.
